% id: 113-05
% book: 베이직쎈 확률과 통계 · 06 이항분포와 정규분포
% section: 유형 067 정규분포곡선의 성질을 이용하여 확률 구하기
% figure: none
% answer: 
% answer_source: 
% variant_level: 0 (원본 전사)
% 이 파일은 build.mjs 가 items.json 에서 만든다. 고칠 때는 items.json 을 고친다.
\smash{\raisebox{1.5mm}{\dmkichul{베이직쎈 확률과 통계 113쪽 05번}}}
정규분포 $\mathrm{N}(m{,}\mkern3mu\mathopen{}\sigma^2)$을 따르는 확률변수 $X$에 대하여\par\smallskip\dispeq{$\mathrm{P}(m-\sigma\le X\le m+\sigma)=a$,}\dispeq{$\mathrm{P}(m-2\sigma\le X\le m+2\sigma)=b$}\par\smallskip\noindent 라 할 때, $\mathrm{P}(m-\sigma\le X\le m+2\sigma)$를 $a$, $b$를 사용하여 나타내면?
\begin{choicesii}{\rule[-3ex]{0pt}{8ex}$\dfrac{a+b}{2}$}{\rule[-3ex]{0pt}{8ex}$\dfrac{a-b}{2}$}{\rule[-3ex]{0pt}{8ex}$\dfrac{b-a}{2}$}{\rule[-3ex]{0pt}{8ex}$\dfrac{2a+b}{2}$}{\rule[-3ex]{0pt}{8ex}$\dfrac{a+2b}{2}$}\end{choicesii}
